\chapter{La Faille de Waterdeep}

\begin{center}
    \textit{Mercredi 15 juillet 2026 --- Session d'initiation}
\end{center}

\vspace{1em}

Kael ouvrit les yeux sur la pierre froide.

Il lui fallut quelques instants pour comprendre où il se trouvait. L'air était humide, chargé de cette odeur particulière propre aux lieux où la lumière du jour ne pénétrait que rarement. Autour de lui, des murs de pierre, une porte fermée et, surtout, aucune explication quant à la manière dont il était arrivé là.

Il n'était cependant pas seul.

Deux autres prisonniers occupaient la même cellule.

Le premier se nommait Morrÿs. Le second, Obélik.

Les présentations furent sommaires. Aucun des trois n'avait vraiment de raison de faire confiance aux deux autres, mais leur situation commune imposait au moins un minimum de coopération. Quelques mots furent échangés, suffisamment pour mettre des noms sur les visages et comparer brièvement leurs situations. Sur un point au moins, tous trois semblaient s'accorder : chacun estimait avoir été jeté en prison à tort, ou tout du moins sans raison qui lui paraisse réellement valable.

Ils ne se connaissaient pas encore, ignoraient s'ils pouvaient se faire confiance, mais partageaient déjà une même certitude : aucun d'entre eux n'avait l'intention de rester enfermé ici.

Avant même de savoir ce qui les attendait au-dehors, leur première aventure venait donc de commencer de la manière la plus simple qui soit :

il fallait sortir de cette cellule.

Ils examinèrent leur prison et cherchèrent un moyen de s'en échapper. Porte, serrure, murs, environnement : tout ce qui pouvait offrir une possibilité de sortie méritait leur attention. Ils n'étaient encore que trois inconnus réunis par les circonstances, mais ils poursuivaient désormais le même objectif.

Et ils finirent par quitter leur cellule.

Ce qui aurait pu n'être qu'une simple évasion prit rapidement une tout autre dimension.

La prison était en plein chaos.

Des kobolds avaient pénétré dans les lieux et leur présence avait bouleversé l'ordre habituel de la forteresse. Les gardes étaient mobilisés, les couloirs n'offraient plus la sécurité que leurs geôliers avaient probablement imaginée, et les trois prisonniers bénéficièrent de cette confusion pour progresser.

Leur fuite les mena jusqu'à une salle des gardes, qui servait de stockage d'armes.

C'est là qu'ils rencontrèrent Alan.

\illustration[0.3\textwidth]{alan.png}
{Alan, tavernier de Waterdeep.}
{alan}

Alan était tavernier à Waterdeep. Comme eux, il avait été enfermé dans cette prison et avait profité du chaos pour s'échapper de sa cellule. C'est au cours de leurs évasions respectives que leurs chemins finirent par se croiser dans cette salle des gardes.

Mais la salle dans laquelle ils venaient d'arriver recelait quelque chose de bien plus inquiétant que quelques armes abandonnées.

Une immense faille s'ouvrait devant eux.

Un gouffre béant dont ils ne pouvaient distinguer le fond déchirait les lieux. Aussi loin que portait leur regard, rien ne permettait d'en mesurer véritablement la profondeur. La pierre disparaissait simplement dans les ténèbres.

Il ne s'agissait plus seulement de quitter une prison. Quelque chose se produisait sous Waterdeep, quelque chose d'assez important pour permettre à des Kobolds d'envahir les lieux.

Kael, Morrÿs et Obélik poursuivirent néanmoins leur progression avec Alan. Chaque salle traversée apportait son lot de dangers et de nouvelles questions.

Puis ils découvrirent la fosse commune.

Une odeur pestilentielle émanait de l'endroit. Kael s'approcha suffisamment pour jeter un œil à l'intérieur. Des corps y avaient été entassés, abandonnés là depuis assez longtemps pour que l'odeur de la chair en décomposition emplisse les lieux.

Lorsqu'il revint vers ses compagnons, Kael leur résuma néanmoins la situation avec un peu moins de gravité : ça sentait la chèvre en décomposition.

Derrière cette comparaison peu flatteuse se cachait une réalité beaucoup plus sombre. La fosse rendait soudain très concrètes les conséquences du chaos qui régnait dans la prison. Ce n'était plus seulement une occasion providentielle de s'évader : des gens étaient morts ici.

Et tous ceux qui vivaient encore n'avaient pas nécessairement l'intention de les laisser partir.

Pour poursuivre leur chemin, les évadés durent encore franchir un dernier obstacle : un immense escalier conduisant au hall de la prison. Une volée interminable de cinquante à cent marches se dressait devant eux.

Mais l'escalier n'avait rien d'ordinaire.

Dès que quelqu'un tentait de s'y engager, la magie de la prison s'activait. Les marches s'aplanissaient brusquement les unes après les autres, transformant l'escalier tout entier en un gigantesque toboggan qui renvoyait l'imprudent vers le bas. Dans le même temps, une sorte de sirène retentissait, une alarme magique annonçant sans aucune discrétion qu'un prisonnier tentait de passer.

Monter ne consistait donc pas simplement à gravir les marches : il fallait trouver un moyen de déjouer le piège avant que celui-ci ne les ramène inlassablement à leur point de départ, tout en signalant leur présence à toute la prison.

\begin{narrateur}
Fort heureusement, nos aventuriers avaient découvert, sur le cadavre encore chaud d'un garde gisant au pied de l'escalier, un code qui aurait sans doute pu leur faciliter considérablement la tâche.

Fort malheureusement, aucun d'entre eux n'eut l'idée de s'en servir.
\end{narrateur}

À la place, ils optèrent pour une méthode autrement plus longue, laborieuse et, à leurs yeux sans doute, parfaitement raisonnable : des pitons et une corde.

Le principe était simple. Gravir quelques marches, planter solidement un piton dans la pierre, y assurer la corde, puis continuer. Lorsque l'alarme magique retentissait et que l'immense escalier se transformait soudain en toboggan, tous s'agrippaient à la corde pour ne pas être précipités jusqu'en bas. Il ne restait alors plus qu'à attendre patiemment que le toboggan redevienne un escalier, gravir quelques marches supplémentaires, planter un nouveau piton, déplacer la corde et recommencer.

Encore.

Et encore.

Marche après marche, piton après piton, ils remontèrent ainsi les cinquante à cent marches de cet interminable escalier.

L'ascension fut longue. Très longue.

Heureusement pour eux, malgré l'alarme magique qui avait consciencieusement annoncé chacune de leurs tentatives, personne ne les rattrapa.

Ils finirent donc par atteindre le sommet, probablement très satisfaits d'avoir triomphé de cet ingénieux système de sécurité sans jamais se demander à quoi pouvait bien servir le code trouvé sur le garde quelques instants auparavant.

Devant eux s'ouvrait enfin le hall de la prison.







Le capitaine des gardes s'y trouvait.

Mais il était loin d'être seul.

Une véritable horde de kobolds avait envahi le hall. Certains combattaient au sol tandis que d'autres, pourvus d'ailes, virevoltaient au-dessus de la mêlée. Au milieu d'eux, le capitaine tenait bon.

Et la différence de niveau était pour le moins évidente.

Là où Kael, Morrÿs, Obélik et Alan avaient parfois besoin de conjuguer leurs efforts pour venir à bout d'un seul kobold, le capitaine en abattait trois, quatre, parfois cinq à la fois.

Arrivés au sommet de leur interminable escalier, nos aventuriers auraient pu profiter de la confusion pour tenter de passer. Ils choisirent pourtant de rejoindre le combat et de prêter main-forte au capitaine.

Le contraste fut quelque peu cruel.

Pendant que celui-ci taillait littéralement son chemin à travers les rangs adverses, les quatre compagnons s'acharnaient à plusieurs sur les créatures qui leur faisaient face. Ils n'étaient peut-être pas encore les héros les plus redoutables de Waterdeep, mais leur aide permit malgré tout de participer à l'élimination des derniers kobolds.

Lorsque le calme revint enfin dans le hall, le capitaine put accorder toute son attention aux nouveaux venus.

La discussion ne tarda pas à devenir délicate.

Il comprit rapidement qu'il avait devant lui quatre des prisonniers qui auraient normalement dû se trouver derrière les barreaux. Les aventuriers tentèrent bien un coup de bluff pour le convaincre du contraire, mais l'expérience ne fut essayée qu'une seule fois : le capitaine n'était manifestement pas dupe.

Pourtant, contre toute attente, il ne chercha pas à les remettre immédiatement en cellule.

La prison n'était plus sûre. Entre les kobolds, les morts et tout ce qui venait de se produire, enfermer de nouveau les quatre évadés n'était manifestement plus sa priorité. Il leur demanda plutôt de le suivre à l'extérieur et leur annonça qu'il les laisserait partir.

Ils sortirent donc avec lui.

À peine dehors, le capitaine eut un pressentiment.

Quelque chose n'allait pas.

Il n'eut pas le temps de réagir.

Une épée surgit de nulle part et le transperça.

\illustration[0.5\textwidth]{capitaine_eventre.png}
{Le capitaine Virgile face à la créature invisible.}
{virgile}

L'attaque avait été aussi brutale qu'inattendue. Aucun adversaire n'était visible, et pourtant le capitaine venait d'être frappé avec une violence suffisante pour l'abattre d'un seul coup.

Les aventuriers parvinrent à lui prodiguer des soins magique et à le ramener à lui.

À peine conscient, il eut juste le temps de leur ordonner de retourner se mettre à l'abri dans le hall.

L'épée invisible frappa de nouveau.

Cette fois, le capitaine parvint, par un heureux hasard, à esquiver le coup.

Le troisième ne lui laissa pas cette chance.

Il s'effondra une nouvelle fois tandis que Kael, Morrÿs, Obélik et Alan se réfugiaient précipitamment à l'intérieur.

Une fois à l'abri dans le hall, un problème assez simple se posa à eux : personne ne savait vraiment quoi faire.

Quelque chose d'invisible et manifestement capable de tuer le capitaine en un seul coup les attendait dehors. Sortir au hasard ne semblait donc pas constituer le meilleur plan qu'ils aient élaboré depuis le début de leur évasion.

C'est alors qu'une petite voix se fit entendre dans leur tête.

Quelqu'un avait une idée.

Une idée risquée.

Et pour qu'elle fonctionne, il fallait que l'un d'entre eux retourne dehors.

Morrÿs accepta.

Il franchit la porte et s'avança à découvert.

Pendant quelques instants, rien ne se passa.

Puis la créature apparut.

Elle ne devint visible que l'espace de quelques millisecondes, juste avant de porter son attaque. Une silhouette se matérialisa devant Morrÿs, une épée déjà levée au-dessus de lui, prête à s'abattre.

Morrÿs n'eut pas le temps de réagir.

La créature non plus.

Une flèche la traversa.

Elle mourut sur le coup.

C'est alors qu'apparut celui qui venait de tirer.

Un halfelin d'une soixantaine d'années s'approcha tranquillement du groupe. Son attitude contrastait singulièrement avec ce qui venait de se produire. Il observa la créature qu'il venait d'abattre d'une seule flèche et résuma la situation avec le plus grand calme :

\begin{quote}
\itshape
« Sale bête. »
\end{quote}

Comme si tuer d'un trait une créature invisible capable de terrasser le capitaine de la prison n'avait rien de particulièrement remarquable.

L'étrange halfelin se porta ensuite au secours du capitaine et lui prodigua les soins nécessaires.

Son attention se tourna alors vers Alan.

Il sembla le reconnaître et lui en fit la remarque.

Alan, lui, nia connaître cet homme.

L'halfelin ne poussa pas davantage la discussion, mais quelque chose dans toute cette affaire semblait le préoccuper. Selon lui, la présence d'une créature aussi puissante à cet endroit n'avait rien de normal. On n'envoyait pas une telle chose ici sans raison.

\illustration[0.5\textwidth]{Sorti_daffaire.png}
{\textit{Morrÿs, Kael, Obélik et leurs étranges compagnons d'évasion.}}
{sorti_d_affaire}

Il conseilla finalement aux quatre aventuriers de quitter Waterdeep et de prendre la direction de Phandalin. Rester dans la cité n'était probablement pas une excellente idée pour trois prisonniers fraîchement évadés. S'ils étaient reconnus, leur liberté nouvellement acquise risquait d'être de courte durée.

Phandalin leur offrirait au moins la possibilité de s'éloigner de la ville et de disparaître quelque temps.

Le conseil donné, l'halfelin aurait parfaitement pu repartir avec eux.

Il choisit une autre direction.

Il retourna vers la prison.

Les mains dans les poches.

Puis, comme si les lois les plus élémentaires de la gravité n'étaient pour lui qu'une vague recommandation, il se mit à marcher sur les murs. Il poursuivit son chemin jusqu'au plafond et s'enfonça ainsi dans la prison, parfaitement à l'aise, avec la ferme intention d'y faire un peu de ménage.

Kael, Morrÿs et Obélik venaient de voir cet étrange vieillard abattre d'une seule flèche une créature qui avait tué deux fois le capitaine des gardes, soigner ce dernier, reconnaître mystérieusement Alan, puis repartir les mains dans les poches en marchant au plafond.

Phandalin semblait soudain être une destination tout à fait raisonnable.